\documentclass[11pt,a4paper,notitlepage,twocolumn]{article}
\usepackage{amsmath,amssymb,amsbsy}
\usepackage{float}
\usepackage[french]{babel}
\usepackage{graphicx}

\usepackage[utf8x]{inputenc}
\usepackage[T1]{fontenc}
\usepackage{palatino}
\usepackage{manfnt}
\usepackage{hyperref}
\usepackage{nicefrac}

\usepackage[active]{srcltx}
\usepackage{scrtime}

\newcommand{\exercice}[1]{\textsc{\textbf{Exercice}} #1}
\newcommand{\question}[1]{\textbf{(#1)}}
\setlength{\parindent}{0cm}

\begin{document}

\title{\textsc{Croissance économique\\ \small{Session de rattrapage}}}
\date{Mardi 16 juin 2026}

\maketitle
\thispagestyle{empty}

Soient $K(t)$ le stock de capital physique d'une économie à l'instant
$t$, $L(t)$ la population (assimilée à la force de travail) à l'instant
$t$ dont le taux de croissance $n>0$ est constant, $T>0$ une quantité
de terre — ou plus généralement de ressources naturelles —
disponible \emph{en quantité fixe}, et $Y(t)$ la production à
l'instant $t$. La production est décrite par la fonction~:
\[
Y(t) = K(t)^{\alpha}L(t)^{\beta}T^{\gamma}
\]
avec $\alpha,\beta,\gamma\in]0,1[$ et $\alpha+\beta+\gamma=1$. La
terre, contrairement au capital et au travail, n'est ni accumulable ni
reproductible~: sa quantité $T$ reste constante au cours du
temps. \textbf{(1)} On suppose que la population à l'instant initial
$t=0$ est connue, $L(0)=L_0$. Déterminer le niveau de la population à
un instant $t$ quelconque. \textbf{(2)} Étudier les rendements
d'échelle de la fonction de production. Montrer en particulier que,
\emph{à quantité de terre donnée}, la production présente des
rendements d'échelle décroissants par rapport au couple
$(K,L)$. Quelle caractéristique de la terre explique ce
résultat~? \textbf{(3)} On note $k(t)=\nicefrac{K(t)}{L(t)}$ le
capital par tête et $z(t)=\nicefrac{T}{L(t)}$ la quantité de terre par
tête. Exprimer la production par tête $y(t)=\nicefrac{Y(t)}{L(t)}$ en
fonction de $k$ et de $z$. \textbf{(4)} À quantité de terre par tête
$z$ donnée, la fonction $k\mapsto f(k)=k^{\alpha}z^{\gamma}$ est-elle
néoclassique~? On vérifiera notamment les conditions
d'Inada. \textbf{(5)} Déterminer la dynamique de la terre par tête
$z(t)$, puis résoudre l'équation différentielle obtenue. Interpréter.
\textbf{(6)} La dynamique du stock de capital agrégé est donnée par~:
\[
\dot K(t) = sY(t)-\delta K(t)
\]
où $s\in]0,1[$ est le taux d'épargne et $\delta>0$ le taux de
dépréciation du capital. Établir l'équation différentielle vérifiée
par le capital par tête $k(t)$. Pourquoi cette équation n'admet-elle
pas d'état stationnaire strictement positif et
constant~? \textbf{(7)} On introduit le ratio capital sur production
$\kappa(t)=\nicefrac{K(t)}{Y(t)}$. En différenciant (le logarithme de)
la fonction de production, montrer que le taux de croissance de la
production vérifie~:
\[
g_Y(t) = \alpha\, g_K(t) + \beta n
\]
\textbf{(8)} En déduire que le ratio capital sur production obéit à
l'équation différentielle linéaire~:
\[
\dot\kappa(t) = (1-\alpha)s - \left[(1-\alpha)\delta + \beta n\right]\kappa(t)
\]
\textbf{(9)} Montrer qu'il existe un unique ratio capital sur
production stationnaire $\kappa^{\star}$ strictement positif, le
calculer, et établir qu'il est globalement stable. Donner l'expression
de $\kappa(t)$. \textbf{(10)} On se place désormais sur le sentier de
croissance équilibrée, c'est-à-dire lorsque $\kappa=\kappa^{\star}$.
Calculer les taux de croissance $g_K$ et $g_Y$, puis celui de la
production par tête $g_y$. Montrer que~:
\[
g_y = -\frac{\gamma n}{1-\alpha}
\]
Interpréter le signe de ce taux de croissance. \textbf{(11)} Discuter
l'effet sur $g_y$ du taux de croissance démographique $n$ et de la
part de la terre $\gamma$ dans la production. Le taux d'épargne $s$
influence-t-il $g_y$~? Que devient le modèle lorsque
$\gamma\rightarrow 0$~? Quel enseignement, déjà présent chez les
économistes classiques (Ricardo, Malthus), retrouve-t-on
ici~? \textbf{(12)} On suppose à présent que la production bénéficie
d'un progrès technique augmentant le travail~:
\[
Y(t) = K(t)^{\alpha}\bigl(A(t)L(t)\bigr)^{\beta}T^{\gamma}
\]
où l'efficacité du travail croît au taux constant
$\nicefrac{\dot A(t)}{A(t)}=x>0$. Reprendre la démarche de la question
\textbf{(10)} et montrer que le taux de croissance de la production
par tête sur le sentier de croissance équilibrée est~:
\[
g_y = \frac{\beta x - \gamma n}{1-\alpha}
\]
Sous quelle condition l'économie connaît-elle une croissance par tête
durablement positive~? Commenter au regard de la question
\textbf{(11)}.

\end{document}
